\section{作业及答案：结构力学基础与技巧班第13次作业及答案（弯矩分配法）}
\begin{analysisbox}
弯矩分配法按固端弯矩、分配系数、传递系数和逐轮平衡进行。最后应检查每个节点的杆端弯矩代数和是否接近外加节点力矩。
\end{analysisbox}
\subsection*{第 1 页转写}
七、练习题

1.基本概念

7-1图7-51中哪一种情况不能用力矩分配法计算。（

）（河海大学2005）

图7-51

7-2

图7-52所示结构AC杆A端的力矩分配系数$\mu$AC=

（浙江大学2009）

7-3图7-53所示结构，要使结点B产生单位转角，则在结点B需施加的外力偶为

）。（合肥工业大学2010、西南交通大学2010）

A.13i

B. 10i

C. 9i

D. 8i

7-4图7-54所示刚架各杆的线刚度=常数，要使A点发生顺时针的单位转角，则

A点需要施加的弯矩m=

。（广西大学2010）

\subsection*{第 2 页转写}
2.5EI

2EI

EI

4m

图7-52

图7-53

图7-54

2.荷载作用下的计算

7-5图7-55所示结构，各杆EI为常数，截面C、D两处的弯矩值Mc、Mp分别为

）。（单位：kN m）（哈尔滨工业大学2009）

A.1.0,2.0

B.2.0,1.0

C.1.0，2.0

D.-2.0，-1.0

7-6

用力矩分配法求图7-56所示结构弯矩图，=1。（广州大学2012）

7-7

用力矩分配法计算图7-57所示结构，并作M图。（计算两轮）（哈尔滨工业大

学2009)

9KN

40kN/m

1.2EI

1.2EI

1.2E

EL

3m

6m

图7-55

图7-56

图7-57

7-8计算图7-58所示结构，并绘M图。要求在力法、位移法、力矩分配法中选择，

EI为常数。（重庆大学2010）

7-9用力矩分配法计算图7-59所示结构，并作M图。已知：q=2.4kN/m，各杆

EI相同。（每结点分配两次）（西南交通大学2008）

FP

EI

2E1

10m

EI

7.5m

10m

4m

7.5m

图7-58

图7-59

\subsection*{第 3 页转写}
7-10用力矩分配法作图7-60所示连续梁的弯矩图。设EI=常数。（山东建筑大学

2012)

30kN-m

7-11用力矩分配法求解图7-61所示结构，

10kN

作M图。EI为常数（忽略剪力和轴力的影响）。（青

岛理工大学2011）

3m

3m

2m

7-12

用力矩分配法计算图7-62所示结构，

图7-60

并作M图，各杆EI相同，为常数。（重庆大学2011)

2kN/m

3kN/m

2kN/m

6kN-m

EI

EI

8kN

4m

4m

4m

3m

4m

图7-61

图7-62

八、练习题答案

7-1

D.

7-2

4/7。提示：C支座相当于固定支座。

7-3

B。提示：BC杆转动刚度为零。

7-4

Ili.

7-5

B.

7-6

M图见图7-63。

M图见图7-64。提示：MA=

7-7

40$\times$32=45kN m。

115.65

68.55

81.14

47.1

0.85qp

0.2qp

0.65qp

0.1gP

23.55

40.57

M图(单位：kN-m)

图7-63

图7-64

7-8M图见图7-65。提示：本题用力矩分配法较方便，C支座相当于固定支座。

7-9M图见图7-66。提示：CFG为静定部分，该部分弯矩图可以直接画出，并且其

在C点的反作用力对左侧超静定部分的弯矩图无影响。

\subsection*{第 4 页转写}
17.36

8.68、

12.08

12.08

10

8.68

45

16.88

15

4.34

6.04

M图($\times$Fp/41)

M图(单位：kN-m)

图7-65

图7-66

7-10

M图见图7-67。

7- 11

M图见图7-68。

12

<8.57

1.71

3.43

20

25

M图(单位：kN-m)

M图（单位：kN-m)

图7-67

图7-68

7-12M图见图7-69。

16

M图(单位：N-m)

图7-69

